% ===================================================================
\chapter{From Plan to Algorithms: The Early Experiments}
\label{ch:journey}

This chapter covers the part of the project that came before any hardware:
what was planned, which experiments were run, what each one showed, and how
two algorithms --- DRHE and a linear predictive coder --- came to be the ones
implemented. It is written in the order things happened. Several conclusions
drawn at the time later turned out to be wrong, and they are reported as they
were believed then, each with a pointer to the chapter where it was
overturned, because the reasons they were wrong are some of this thesis's
results.

\section{The plan}
\label{sec:jr_plan}

The Thesis~A report (April 2026) set three goals, each aimed at a limitation
it attributed to Kiem's work.

\begin{enumerate}
\item \textbf{Mitigate quantisation error.} Kiem converts the range-FFT output
to 16-bit fixed point (FX16) before compressing, and that step is where his
detection errors come from (Table~\ref{tab:kiem_t35}: FX16 loses
\SI{0.59}{\percent} of detections relative to FP16). Two remedies were
proposed: a non-linear, piecewise companding transform applied before
quantisation, borrowed from synthetic aperture radar image compression
\cite{deng2024}; or a floating-point implementation, justified by the large
latency headroom in Kiem's design.
\item \textbf{Evaluate other compression algorithms}, in particular NXP's
CP4/CP8 block-floating-point formats, Block Adaptive Quantization and the
learned spectral scheme AdaRadar, against DRHE.
\item \textbf{Run the whole chain end to end on hardware}, with an evaluation
module streaming into an FPGA compressor and a PC decompressing.
\end{enumerate}

\noindent
Goal~1 was investigated and turned into a finding about compression rather
than a remedy for quantisation (Section~\ref{sec:jr_quant}). Goal~2 was carried
out with a partly different list of algorithms (Section~\ref{sec:jr_bench}).
Goal~3 was carried as far as possible without a board: implementation through
placement and routing, but no silicon.

\begin{figure}[htbp]
\centering
\begin{tikzpicture}[font=\scriptsize,x=1cm,y=1cm]
\draw[slate,line width=0.8pt,-{Latex[length=2mm]}] (0,0) -- (15.0,0);
\foreach \x/\l/\d/\s in {
  0.6/{Thesis A plan;\\Kiem's workflow\\rebuilt in MATLAB}/{Mar--Apr}/1,
  2.4/{synthetic scenes:\\DRHE CR $\approx$ 1.1}/{Mar--May}/-1,
  4.2/{Ancortek capture\\attempt; move to\\ColoRadar}/{by Jul}/1,
  6.0/{quantisation (LFR)\\study; 8-algorithm\\benchmark}/{to Jul}/-1,
  7.8/{HLS: DRHE and LPC,\\Phases 1--3,\\two defects}/{to 15 Sep}/1,
  9.6/{audit: predictors\\not helping;\\TDM axis fix}/{15 Sep}/-1,
  11.4/{Vivado, replica of\\Kiem, ablation\\ladder}/{23--24 Sep}/1,
  13.4/{replication of\\Kiem's simulation;\\re-verification}/{24 Sep}/-1}{
  \fill[accent] (\x,0) circle (2.2pt);
  \draw[slate!60] (\x,0) -- (\x,0.35*\s);
  \ifnum\s=1
    \node[align=center,anchor=south,text=slate] at (\x,0.35) {\textcolor{accent}{\textbf{\d}}\\\l};
  \else
    \node[align=center,anchor=north,text=slate] at (\x,-0.35) {\textcolor{accent}{\textbf{\d}}\\\l};
  \fi
}
\end{tikzpicture}
\caption[Project timeline]{The order in which the work was done. Dates are
from the supervisory meeting log and the repository history; the chapters of
this report follow the same order, with the replication and comparison with
Kiem placed after the hardware results.}
\label{fig:timeline}
\end{figure}

\section{Rebuilding Kiem's workflow}
\label{sec:jr_rebuild}

The first piece of work was a MATLAB simulation environment,
\texttt{Matlab\_Sim}, modelled stage for stage on Kiem's workflow
(Section~\ref{sec:kiem_workflow}): synthetic or recorded ADC input, shift to
32-bit full scale, Hann window, range FFT divided by $N$, first half kept,
compression and decompression, Doppler FFT, NCI over all receive channels,
peak detection and a comparator computing CR, FN\%, FP\%, noise floor and SNR.
Where Kiem does not specify a step, a choice was made and recorded; two of those
choices matter later.

\paragraph{The FX16 scale.} \texttt{compress\_fx16.m} maps a full-scale ADC
sinusoid, after windowing and the FFT, to the largest \texttt{int16} value,
Equation~\eqref{eq:scale}. Kiem's hardware instead keeps the top 16 bits of the
32-bit FFT word, in which a full-scale sinusoid is only about 8191. The two
scalings differ by a factor of four --- two bits --- and
Section~\ref{sec:kr_early} shows that this difference alone moves a white-noise
compression ratio substantially.

\paragraph{The noise-floor reference.} 0\,dBFS was taken as the NCI power of
a full-scale sinusoid carried through the whole chain
(\texttt{computeFullScaleRefPower.m}). Kiem does not say what his reference is.

\section{Synthetic scenes, and a compression ratio of 1.1}
\label{sec:jr_synth}

Before any real data was available, the environment was exercised with
synthetic scenes built to resemble Kiem's: 512 samples, 512 ramps, four
receive channels, a 14-bit ADC, some two hundred point targets to give about
his 170 detections, and additive white Gaussian noise set so that the measured
noise floor and SNR matched his Table~3.5. The target was his DRHE CR of 3.25.
Table~\ref{tab:jr_synth} shows what came out.

\begin{table}[htbp]
\centering
\caption[Early synthetic-scene results]{DRHE on the early synthetic scenes
(\texttt{main\_final.log}, \texttt{kiem\_final\_with\_noise.log},
\texttt{kiem\_enob\_sweep.txt}). ``ENOB'' is an optional reduction of the FX16
code to fewer effective bits. Every scene uses the Matlab\_Sim FX16 scaling.}
\label{tab:jr_synth}
\small
\resizebox{\textwidth}{!}{%
\begin{tabular}{@{}lrrrrrr@{}}
\toprule
Scene & ADC $\sigma$ (codes) & ref.\ det. & NF (dBFS) & SNR (dB) & FN\,/\,FP (\%) & DRHE CR\\
\midrule
Matched to Kiem's NF (\texttt{match-nf})       & 418.2 & 160 & $-71.03$ & 23.81 & 0 / 0 & 1.11\\
Matched to Kiem's FN/FP (\texttt{match-fnfp})  & 418.2 & 160 & $-70.70$ & 23.51 & 0.62 / 0 & 1.10\\
Dense targets, white floor (\texttt{clean})    & 591.1 & 179 & $-71.01$ & 23.23 & 1.68 / 0 & 1.12\\
Mixed white + correlated noise                 & 497.0 & 167 & $-70.76$ & 23.54 & 0 / 0 & 1.09\\
Noise lowered 14\,dB (\texttt{match-cr})       & 78.8  & 163 & $-85.37$ & 38.05 & 0 / 0 & 1.56\\
Noise lowered 24\,dB (\texttt{aggressive})     & 24.9  & 163 & $-95.08$ & 47.76 & 0 / 0 & 2.08\\
Strong zero-Doppler clutter (\texttt{clutter}) & 12.5  & 51  & $-101.99$ & 61.56 & 0 / 0 & 2.57\\
\midrule
Kiem, Table 3.5 (real data)                    & ---   & $\approx$170 & $-70.69$ & 22.33 & 0.59 / 1.18 & 3.25\\
\bottomrule
\end{tabular}}
\end{table}

\noindent
Every scene that matched Kiem's noise floor compressed to about 1.1. Lowering
the noise raised the ratio, but even 24\,dB less noise reached only 2.08, and
then the noise floor and SNR no longer resembled his. A further sweep
(\texttt{noise\_type\_sweep.txt}) replaced white noise with slow-time-correlated
noise: first-order autoregressive noise with correlation up to 0.999, coloured
noise, $1/f$ noise and mixtures. None exceeded CR 1.8, and the correlated
variants that raised CR did so while destroying detection (FN up to
\SI{67}{\percent}), because strongly correlated noise piles up at low Doppler
and is detected as targets.

\paragraph{The interpretation at the time.} The executive summary explained the
gap as a property of white noise: ``simulated AWGN data lacks correlation
across consecutive chirps in the slow-time (Doppler) dimension. Because
algorithms like DRHE rely on predicting sample values from previous ramps, the
lack of correlation produced large residual values.'' Real data, the
reasoning continued, compresses because its floor contains coherent clutter
that DRHE can predict. That explanation motivated the move to real data.

\begin{keybox}
\textbf{What this turned out to be.} The explanation was wrong in its
mechanism. On a white-noise scene the compression ratio is set almost entirely
by how many FX16 least significant bits the noise occupies, and prediction
contributes little either way: on Kiem's own simulated scene, coding the FX16
values with \emph{no prediction at all} gives a ratio within about ten per cent
of DRHE's at every noise level, and often a better one
(Section~\ref{sec:kr_sweep}). The early scenes
compressed to 1.1 because matching Kiem's stated noise floor with white noise
and the Matlab\_Sim scaling puts about 130 LSB of noise into every FX16 value,
roughly thirty times what Kiem's own CR implies. Chapter~\ref{ch:replication}
rebuilds his simulation independently and reproduces the 1.1 from first
principles, and Section~\ref{sec:decomposition} shows that the real data's
higher ratio did not come from predictable clutter either.
\end{keybox}

\section{Looking for real data}
\label{sec:jr_data}

The first attempt at real data used a laboratory Ancortek FMCW module. MATLAB
scripts were written to control it and read its captures
(\texttt{readAncortekBin.m}, \texttt{getAncortekChannels.m}), and several
single-target and background captures were recorded. They were not usable for
this study: the interface scripts were unreliable, and the module's bandwidth
gave too coarse a range resolution to produce the dense range profiles a
compression study needs.

The project therefore moved to the public ColoRadar dataset \cite{kramer2022},
which contains raw ADC captures from exactly the cascade described in
Section~\ref{sec:sensor}: four AWR2243 devices, 12 transmitters and 16
receivers in time-division MIMO. A loader (\texttt{loadColoRadarFrame.m})
reads a frame, applies the TDM flattening of Equation~\eqref{eq:rampindex} and
removes the per-ramp DC term. Fifty consecutive frames of one indoor sequence
were used for everything that follows. This choice also changed the problem in
a way nobody noticed for months: a single-transmitter algorithm was about to be
applied to a twelve-transmitter radar.

\section{The quantisation question}
\label{sec:jr_quant}

Goal~1 was pursued with a study (\texttt{compare\_quantisation\_lfr.m}) on one
ColoRadar frame. Four quantisers were compared, each followed by the same
lossless DRHE coder:

\begin{itemize}
\item \textbf{FP16}: half-precision floating point, the fidelity reference;
\item \textbf{FX16-uniform}: the pipeline's fixed full-scale scaling;
\item \textbf{FX16-linear}: the same uniform quantiser, but scaled so the
frame's own peak fills the \texttt{int16} range;
\item \textbf{FX16-LFR}: the two-stage low-frequency-suppression companding
curve of Deng and Huang \cite{deng2024} --- a linear stretch of small magnitudes
and logarithmic compression of large ones --- with its parameters fitted to the
frame's magnitude distribution.
\end{itemize}

\begin{table}[htbp]
\centering
\caption[Quantisation study]{The quantisation study, ColoRadar frame 0, all
16 receive channels (re-run 24 September 2026, \texttt{\_verify/lfr\_rerun.log}).
``relRMS'' is the RMS error of the range-FFT values relative to their RMS;
``RD error'' is the RMS dB error of the range--Doppler map; ``int16 use'' is the
largest code as a fraction of 32767.}
\label{tab:jr_lfr}
\small
\begin{tabular}{@{}lrrrrr@{}}
\toprule
Quantiser & relRMS (\%) & mag.\ PSNR (dB) & RD error (dB) & int16 use (\%) & DRHE CR\\
\midrule
FP16 (reference)   & 0.0205 & 101.97 & 0.0049 & ---   & ---\\
FX16-uniform       & 1.8636 & 64.44  & 0.4336 & 1.33  & \textbf{3.320}\\
FX16-linear        & 0.0278 & 101.09 & 0.0054 & 89.20 & 1.038\\
FX16-LFR           & 0.0048 & 114.38 & 0.0011 & 94.28 & 0.708\\
\bottomrule
\end{tabular}
\end{table}

Table~\ref{tab:jr_lfr} is one of the most informative results of the whole
project, although it was not read that way at the time. The uniform FX16 used by
Kiem and by this pipeline uses \SI{1.33}{\percent} of the \texttt{int16}
range: the largest code in the frame is 436. Its quantisation error is
correspondingly large. Simply scaling the data to fill the container cuts the
error by a factor of 67 and brings the fixed-point path to within a fraction of
a decibel of FP16. Companding does better still.

But the compression ratio moves the other way. With the container filled,
lossless DRHE compresses to 1.04; with companding, which spreads the values
evenly over the whole code range, it \emph{expands} the data. The quantisation
error and the compression ratio are two sides of the same unused bits. The
uniform code leaves about six of its sixteen bits permanently zero at the top
and rounds away everything below its coarse step. Stretching the data to fill
the container turns those six bits into fine detail at the bottom --- detail
that is mostly receiver noise, which a lossless coder must keep, and which is
exactly what the uniform code's entropy coder was never asked to pay for.

\paragraph{What was decided.} Goal~1, as stated, could be met --- the study
shows the quantisation error can be removed --- but only by giving up almost all
of the compression the rest of the project was about. The pipeline therefore
kept Kiem's fixed FX16 scaling, so that results stayed comparable with his, and
the trade-off was recorded rather than resolved. Section~\ref{sec:decomposition}
later measured it precisely: \num{1.64} of the no-prediction ratio of
\num{3.489} comes from unused container bits.

\paragraph{Floating point instead.} The other remedy, floating point, was
tested in the same spirit: DRHE was run on FP16 values instead of FX16 codes
(``DRHEfp16'', \texttt{compress\_drhe\_fp16.m}). It compresses to 1.04
(Table~\ref{tab:phase1}), because a half-precision mantissa keeps fine noise
detail at every magnitude and the differences never settle near zero. The
conclusion was that the compressor's \emph{input} must be integer. The
\emph{arithmetic inside} the predictor is a separate question: Kiem's own future
work suggested floating point there (Section~\ref{sec:kiem_open}), and the
hardware DRHE of this thesis uses it (Section~\ref{sec:drheimpl}).
Chapter~\ref{ch:interpret} measures what that costs.

\section{Choosing what to compare}
\label{sec:jr_choose}

Goal~2 was narrowed to schemes that could be implemented and scored in the same
MATLAB chain. Each works on the range-FFT output.

\begin{description}[leftmargin=!,labelwidth=26mm,style=nextline,font=\normalfont\bfseries]
\item[FP16, FX16] The two reference formats, as in Kiem.
\item[DRHE] Kiem's final scheme (Section~\ref{sec:kiem_final}):
\texttt{compress\_drhe.m}. Section~\ref{sec:kc_audit} shows it is Kiem's
Section~3.5.3 exactly.
\item[DRHEfp16] The same predictor on FP16 values, to test the floating-point
remedy of Goal~1.
\item[BAQ] Block adaptive quantisation, implemented as a block-floating-point
quantiser: blocks of 16 range bins share an 8-bit exponent chosen from the
block maximum, and each real and imaginary part keeps a 4-bit mantissa. The CR
is fixed by construction at $32/(8 + 8/16) = 3.76$. It stands in for the
block-floating-point family (NXP CP4/CP8 and Kiem's CMPRA/CMPRB are of the same
kind) and for BAQ proper, whose statistics-driven step size
\cite{kwok1989} it approximates with a maximum-driven one.
\item[RDVLE] Range-dependent variable-length encoding from Li's thesis
\cite{li2014}: fewer bits for far range bins, following the $1/R^4$ fall of
received power, here capped at three bits removed.
\item[RDVLE+UDRE] RDVLE followed by Li's uniform dynamic-range re-encoding of
blocks of 16 bins at six bits.
\item[LPC-Huffman] A linear predictive coder, designed for this thesis and
described next.
\end{description}

\noindent
NXP's formats were not implemented separately, since their exact bit layouts
are proprietary and the block-floating-point quantiser represents the family.
AdaRadar \cite{park2026} was left out of scope once the thesis narrowed to
lossless compression; it is a lossy, learned scheme whose decoder runs on the
central processor.

\section{Where this thesis's LPC came from}
\label{sec:jr_lpc}

The name was taken from Kiem's Section~2.7.4. The design was not. Kiem
describes LPC-Huffman as he understood \cite{meucci2022}: a $p$-th order linear predictor
along \emph{fast time}, predicting each raw ADC sample from the previous $p$
samples of the same chirp, with least-squares coefficients and a Huffman code
built from the data. He excluded it without running it
(Section~\ref{sec:kiem_space}).

Taken literally, that method would have been compared with DRHE across four
differences at once: a different domain (raw ADC against range FFT), a
different axis (fast time against slow time), a different coder (data-derived
Huffman against a fixed dictionary) and a different predictor. Any difference in
CR could have come from any of them. The LPC built here was therefore designed
as a \emph{controlled} comparison, changing only the predictor:

\begin{itemize}
\item the same input as DRHE: FX16 range-FFT codes;
\item the same axis: slow time, predicting ramp $m$ of a range bin from earlier
ramps of that bin;
\item the same entropy coder: S4 category, APPEND bits, and Kiem's fixed
Huffman dictionary;
\item a different predictor: a second-order autoregressive model
$\hat{x}[m] = a_1 x[m-1] + a_2 x[m-2]$ applied separately to the real and
imaginary parts, with Yule--Walker coefficients fitted per (range bin, receive
channel, part) over the whole frame and sent as 16-bit side information.
\end{itemize}

\noindent
The idea is simple. DRHE assumes a shape --- constant magnitude, constant phase
increment --- and hard-wires filters for it. The LPC assumes only that each
bin's slow-time sequence is roughly a second-order autoregressive process, and
lets the data choose the coefficients. If DRHE's assumption is right, the fitted
model should find something close to it; if it is wrong, the fitted model should
find something better, or at least find that there is little to predict.
Table~\ref{tab:lpc_vs_kiemlpc} sets the two methods called ``LPC'' side by side.

\begin{table}[htbp]
\centering
\caption[Two methods called LPC]{Kiem's LPC-Huffman (his Section~2.7.4,
after \cite{meucci2022}) and the LPC of this thesis. They share a name and the idea of
least-squares linear prediction, and nothing else that affects the result.}
\label{tab:lpc_vs_kiemlpc}
\small
\begin{tabular}{@{}p{33mm}p{50mm}p{50mm}@{}}
\toprule
& \textbf{LPC-Huffman in Kiem (after \cite{meucci2022})} & \textbf{LPC in this thesis}\\
\midrule
Input & raw ADC samples (assumed by Kiem) & FX16 range-FFT codes\\
Prediction axis & fast time: previous $p$ samples of the same chirp & slow time: ramps $m-1, m-2$ (later $m-n_{\mathrm{Tx}}, m-2n_{\mathrm{Tx}}$) of the same range bin\\
Order & $p$ (not stated) & 2\\
Coefficients & least squares per signal & Yule--Walker per (bin, RX, part), 16-bit\\
Entropy coder & Huffman on residual values, dictionary from the data & S4 + APPEND, Kiem's fixed S4 dictionary\\
Streaming & yes, per chirp & no: needs the whole frame\\
What it is a test of & compressing the ADC stream & DRHE's fixed predictor against a fitted one\\
Evaluated by Kiem & no (CR 2.056 quoted from \cite{meucci2022}) & ---\\
Evaluated here & Chapter~\ref{ch:replication} & Chapters~\ref{ch:results} and~\ref{ch:replication}\\
\bottomrule
\end{tabular}
\end{table}

\section{The benchmark}
\label{sec:jr_bench}

All eight schemes were run over the same 50 consecutive ColoRadar frames with
all 16 receive channels (\texttt{main\_simulation\_coloradar\_batch.m}). The
reference is the double-precision path; FP16 is scored against it like every
other scheme. Table~\ref{tab:phase1} gives the result. It was first produced during the
Thesis~A term, published in the executive summary, and re-run for this report on
24 September 2026 with identical results in every cell
(\texttt{\_verify/phase1\_rerun.log}).

\begin{table}[htbp]
\centering
\caption[The eight-scheme benchmark]{The eight-scheme benchmark: 50-frame
average, 16 receive channels, \num{9934} reference detections. FP\% is a false
discovery rate (Section~\ref{sec:detectionmetrics}). NF is the noise floor.}
\label{tab:phase1}
\small
\begin{tabular}{@{}lrrrrr@{}}
\toprule
Algorithm & CR & FN\% & FP\% & NF (dBFS) & SNR (dB)\\
\midrule
FP16          & 1.00 & 0.00  & 0.00 & $-105.638$ & 35.504\\
DRHEfp16      & 1.04 & 0.05  & 0.01 & $-105.637$ & 35.505\\
FX16          & 1.00 & 3.77  & 1.11 & $-105.374$ & 35.359\\
\textbf{DRHE} & \textbf{3.30} & 3.77 & 1.11 & $-105.374$ & 35.359\\
BAQ           & 3.76 & 10.44 & 9.69 & $-105.080$ & 35.081\\
RDVLE         & 1.12 & 24.01 & 4.62 & $-102.275$ & 33.165\\
RDVLE\_UDRE   & 2.56 & 28.73 & 4.44 & $-101.782$ & 33.004\\
\textbf{LPC}  & \textbf{3.37} & 3.77 & 1.11 & $-105.374$ & 35.359\\
\bottomrule
\end{tabular}
\end{table}

Four observations decided what happened next.

\paragraph{The lossless coders share FX16's detection metrics exactly.}
FX16, DRHE and LPC report identical FN\%, FP\%, noise floor and SNR. This is not
a coincidence but a tautology: both coders reconstruct the FX16 codes
bit-exactly, so the detector sees identical data. What those rows show is the
cost of \emph{quantisation}: \SI{3.77}{\percent} of detections missed and
\SI{1.11}{\percent} of reported detections spurious.

\paragraph{The lossy schemes pay for their ratio in detections.}
BAQ reaches a nominally better ratio --- \num{3.76} against \num{3.30} --- but
misses \SI{10.44}{\percent} of detections instead of \SI{3.77}{\percent}, and
\SI{9.69}{\percent} of what it reports is spurious against \SI{1.11}{\percent}.
That is nearly three times the miss rate and nearly nine times the false
discovery rate for a \SI{14}{\percent} better ratio. RDVLE and RDVLE+UDRE are
worse on both counts. In a system that brakes a car, none of these trades is
acceptable, which is the empirical form of the argument of
Section~\ref{sec:lossless}.

\paragraph{Floating-point storage does not compress.} DRHEfp16 reaches 1.04,
confirming the result of Section~\ref{sec:jr_quant}.

\paragraph{Two predictive coders reach the target.} DRHE at 3.30 and LPC at
3.37 both exceed the 2.4 a Gigabit Ethernet link would need, losslessly.

\section{The decision, and what was expected}
\label{sec:jr_decision}

DRHE and LPC were both taken forward to hardware. The reasoning recorded in the
executive summary was that LPC compresses slightly better (3.37 against 3.30)
but ``requires more complex arithmetic operations on an FPGA, such as
division'', while DRHE ``requires only simple arithmetic logic''; implementing
both would show where the trade really lay.

\begin{keybox}
\textbf{Five beliefs held at this point, and what became of them.}
\begin{enumerate}[leftmargin=*]
\item \emph{DRHE's 3.30 confirmed the implementation, because it matched
Kiem's 3.25.} The match is a coincidence of two different mechanisms on two
different datasets (Chapters~\ref{ch:replication} and~\ref{ch:kiemcompare}).
\item \emph{Both predictors were exploiting Doppler redundancy.} Neither was:
as implemented, each did worse than no prediction (Section~\ref{sec:decomposition}),
because the radar is twelve-transmitter TDM (Section~\ref{sec:predaxisresults}).
\item \emph{LPC compresses better because a fitted model beats a fixed one.}
True for a different reason than expected: on the wrong axis, LPC's fitted
coefficients shrink towards zero so it nearly switches itself off, while DRHE's
fixed filters cannot (Section~\ref{sec:lpcshrink}).
\item \emph{LPC is the expensive design because it divides.} Wrong: the
division runs 1024 times per frame; LPC's cost is its frame buffer and its
latency (Section~\ref{sec:phase3}).
\item \emph{Real data compresses because its noise floor contains predictable
clutter.} Wrong: \num{1.64} of the ratio is unused container bits and
\num{2.13} spectral sparsity (Section~\ref{sec:decomposition}).
\end{enumerate}
\end{keybox}
